\chapter{Statistical Mechanics Background}
\label{ch:statmech}

Every engine in \qv{} is a Markov-chain Monte Carlo (MCMC) sampler of a
Gibbs--Boltzmann distribution, driven slowly toward the zero-temperature (or
zero-quantum-fluctuation) limit. This chapter builds that machinery from
scratch: the Boltzmann distribution, why sampling it at low temperature
solves optimisation problems, Markov chains, detailed balance, and the
Metropolis rule. Readers fluent in MCMC can skim to
Chapter~\ref{ch:sa}.

\section{The Boltzmann distribution}

A system with configuration space $\Omega$ (for us,
$\Omega=\{-1,+1\}^n$) and energy function $E:\Omega\to\mathbb{R}$, held in
equilibrium at temperature $T$, occupies configuration $s$ with probability

\begin{keyeq}[Gibbs--Boltzmann distribution]
\begin{equation}
\pi_\beta(s)=\frac{\mathrm{e}^{-\beta E(s)}}{Z(\beta)},
\qquad
Z(\beta)=\sum_{s\in\Omega}\mathrm{e}^{-\beta E(s)},
\qquad
\beta=\frac{1}{k_BT}.
\label{eq:boltzmann}
\end{equation}
\end{keyeq}

$\beta$ is the \emph{inverse temperature}; in \qv{} it is always
dimensionless (energies carry the units). $Z(\beta)$ is the \emph{partition
function}. Two limits explain the whole strategy of annealing:

\begin{itemize}
\item $\beta\to0$ (infinite temperature): $\pi_\beta$ is uniform over all
$2^n$ configurations --- pure exploration, no preference for low energy.
\item $\beta\to\infty$ (zero temperature): the ratio
$\pi_\beta(s)/\pi_\beta(s^\ast)=\mathrm{e}^{-\beta\,(E(s)-E(s^\ast))}\to0$
for every non-optimal $s$, so \emph{all probability mass concentrates on the
ground states}. Sampling at large $\beta$ \emph{is} optimisation.
\end{itemize}

The catch is that sampling $\pi_\beta$ directly at large $\beta$ is as hard
as the optimisation problem itself: naive samplers get trapped in local
minima because the distribution is a set of sharp, well-separated peaks.
Annealing resolves this by \emph{starting} at small $\beta$, where the
distribution is easy to sample, and increasing $\beta$ slowly enough that the
sampler tracks the shifting distribution.

\section{Markov chains and stationarity}

We generate samples by a Markov chain: a random walk on $\Omega$ whose next
state depends only on the current one, via a transition kernel
$P(s\to s')\ge0$ with $\sum_{s'}P(s\to s')=1$. A distribution $\pi$ is
\emph{stationary} for $P$ if it is preserved by one step of the chain:
\begin{equation}
\sum_{s}\pi(s)\,P(s\to s')=\pi(s')\qquad\text{for all } s'.
\label{eq:stationary}
\end{equation}
If, additionally, the chain is \emph{irreducible} (any state reachable from
any other) and \emph{aperiodic}, the ergodic theorem guarantees that the
chain's state distribution converges to $\pi$ from any start, and long-run
averages along the chain converge to expectations under $\pi$.

\subsection{Detailed balance}

Verifying Eq.~\eqref{eq:stationary} directly is awkward. A stronger, local
condition that implies it is \emph{detailed balance}:

\begin{keyeq}[Detailed balance]
\begin{equation}
\pi(s)\,P(s\to s') \;=\; \pi(s')\,P(s'\to s)
\qquad\text{for every pair } s,s' .
\label{eq:detailed-balance}
\end{equation}
\end{keyeq}

Summing both sides over $s$ immediately recovers
Eq.~\eqref{eq:stationary}. Detailed balance says the equilibrium probability
flux between any two states is symmetric --- the chain, run at equilibrium,
is statistically indistinguishable from its time reversal.

\section{The Metropolis rule, derived}
\label{sec:metropolis}

Split each transition into a \emph{proposal} $g(s\to s')$ (in \qv: pick a
spin uniformly and flip it, so $g$ is symmetric, $g(s\to s')=g(s'\to s)$) and
an \emph{acceptance probability} $A(s\to s')$, so
$P(s\to s')=g(s\to s')\,A(s\to s')$ for $s'\ne s$. Substituting into
Eq.~\eqref{eq:detailed-balance} with $\pi=\pi_\beta$ from
Eq.~\eqref{eq:boltzmann}, the proposal factors cancel and the partition
function cancels too:
\begin{equation}
\frac{A(s\to s')}{A(s'\to s)}
=\frac{\pi_\beta(s')}{\pi_\beta(s)}
=\mathrm{e}^{-\beta\,\Delta E},
\qquad \Delta E = E(s')-E(s).
\end{equation}
Any $A$ satisfying this ratio works; the choice that maximises acceptance
(and hence mixing speed) is the \emph{Metropolis--Hastings} rule:

\begin{keyeq}[Metropolis acceptance]
\begin{equation}
A(s\to s')=\min\!\bigl(1,\;\mathrm{e}^{-\beta\,\Delta E}\bigr)
=\begin{cases}
1 & \Delta E\le 0 \quad\text{(downhill: always accept)}\\[2pt]
\mathrm{e}^{-\beta\,\Delta E} & \Delta E>0 \quad\text{(uphill: accept with Boltzmann weight)} .
\end{cases}
\label{eq:metropolis}
\end{equation}
\end{keyeq}

This is exactly what the C++ core executes; for a single-spin flip the
$\Delta E$ comes from the cached-local-field formula Eq.~\eqref{eq:deltaE},
making each proposal $O(1)$.

\begin{physbox}
At small $\beta$ almost every move is accepted and the walker diffuses freely
across the landscape. At large $\beta$ uphill moves are exponentially
suppressed and the walker rolls downhill into the nearest basin. Annealing
is the art of spending long enough at intermediate $\beta$ --- where the
walker can still cross barriers whose height is a few times $1/\beta$ ---
that it ends in the \emph{right} basin before freezing.
\end{physbox}

\subsection{Sweeps, ergodicity, and shuffled visit order}

A \emph{sweep} is $n$ proposals --- on average one per spin. \qv{} visits
spins in a random order re-shuffled each sweep (rather than uniformly at
random with replacement), which lowers the chance of neglecting a spin for
long stretches while preserving detailed balance. All engine parameters
counting ``sweeps'' (\code{sweeps\_per\_beta}, \code{sweeps\_per\_step},
\code{worldline\_sweeps}, \code{cluster\_sweeps}) are in units of full
sweeps.

\section{From sampling to annealing}
\label{sec:annealing-idea}

Simulated annealing~\cite{kirkpatrick} runs Metropolis dynamics while
increasing $\beta$ along a \emph{schedule}
$\beta_1<\beta_2<\dots<\beta_K$, spending a fixed number of sweeps at each
value. Two classical facts frame all schedule design:

\begin{enumerate}
\item \textbf{(Geman--Geman)} With a logarithmically slow schedule
$\beta_k \propto \log k$, SA converges to a global optimum with probability
one. This is far too slow to use, but it shows the trade-off is fundamental:
faster cooling risks freezing into excited states.
\item \textbf{(Acceptance targeting)} A schedule is well-matched to a problem
when the \emph{uphill acceptance rate} decays from near~1 at the start to
near~0 at the end. Given a characteristic uphill move size $\Delta$, the
acceptance at inverse temperature $\beta$ is
$p=\mathrm{e}^{-\beta\Delta}$, hence
\begin{equation}
\beta(p)=\frac{-\ln p}{\Delta}.
\label{eq:beta-from-p}
\end{equation}
\qv's auto-tuned schedules (Section~\ref{sec:tuned-schedules}) measure
$\Delta$ from the problem itself --- the 75th percentile of $|\Delta E|$ over
random single-spin flips --- and place $\beta_{\text{start}},
\beta_{\text{end}}$ at target acceptances via Eq.~\eqref{eq:beta-from-p}.
No manual temperature tuning is required.
\end{enumerate}

\begin{notebox}[Quantum analogue]
Chapters~\ref{ch:sqa}--\ref{ch:ctpimc} replace ``thermal fluctuations
controlled by $1/\beta$'' with ``quantum fluctuations controlled by a
transverse field $\Gamma$''. The mathematical machinery of this chapter
survives intact: the quantum partition function is mapped (exactly in a
well-defined limit) onto a \emph{classical} Boltzmann distribution in one
extra dimension, and the very same Metropolis rule
Eq.~\eqref{eq:metropolis} samples it.
\end{notebox}
